\section{Background and Related Work}
\label{sec:background}

\subsection{Kubernetes Resource Semantics}

A container may declare CPU and memory requests, limits, both, or neither.
Requests form the scheduler's placement baseline: a Pod is feasible on a Node
only when its effective requests fit the Node's remaining allocatable
resources. Limits instead define enforcement ceilings used by the kubelet and
container runtime \citep{k8s_resource_management}. A namespace
\textit{LimitRange} may insert defaults during admission, and those effective
values are visible in the stored Pod specification
\citep{k8s_limit_range}.

Requests and limits also contribute to Pod Quality of Service classification,
but QoS is not a direct proxy for current work. A \textit{Burstable} Pod can
legitimately consume far more CPU than its request when its limit permits it,
while a Pod with a large request can be idle
\citep{k8s_pod_qos}. Runtime protection therefore requires an observation of
current consumption as well as a baseline against which to interpret it.

Metrics Server exposes recent CPU and memory usage through the Kubernetes
Metrics API \citep{k8s_metrics_server}. CPU usage is represented in cores or
millicores and memory in bytes. These observations are point-in-time,
best-effort signals rather than reservations, so unavailable or incomplete
measurements require an explicit safety policy.

\subsection{Descheduling and Runtime Awareness}

The Descheduler runs after initial placement and proposes Pod evictions
according to a configured strategy. Before the eviction request reaches the
Kubernetes API, evictor plugins can enforce general eligibility and
pre-eviction checks. Multiple pre-eviction filters may be chained, so every
enabled filter must permit the candidate
\citep{k8s_descheduler_docs}. This extension point allows runtime protection
to remain independent of candidate ordering and the strategy's placement
objective.

Runtime telemetry has been used to revise Kubernetes placement at several
decision stages. The closest work combines load-aware scheduler scoring with a
custom Descheduler evictor that uses CPU, memory, network, and disk telemetry
to select the Pod with the greatest positive placement-score improvement
\citep{marchese2025enhancing}. An earlier network-aware scheduler and
Descheduler operator similarly uses measured latency and inter-service traffic
to evict a Pod when another Node offers a better score
\citep{marchese2022network}. Kinitos uses policy-specific scheduler and
Descheduler plugins to identify violated network dependencies and trigger
rescheduling, retaining the current placement when no feasible destination
exists \citep{tsokov2025kinitos}. Outside the Descheduler framework, Nautilus
uses runtime QoS and resource pressure to select and migrate a microservice
through a scale-out-before-delete workflow \citep{fu2022adaptive}. These
systems use telemetry for scoring, candidate selection, destination selection,
or managed migration. The Actual Usage Evictor operates at a narrower stage:
after a host strategy selects a candidate and immediately before eviction, its
global \textit{PreEvictionFilter} only permits or blocks that candidate. It
does not select, rank, place, or migrate Pods; it compares live CPU and memory
with effective requests and fails closed when metrics are unavailable or
incomplete.

The objective also differs from kubelet node-pressure eviction. Under resource
pressure, the kubelet ranks Pods to reclaim a scarce Node resource using Pod
Priority, whether usage exceeds requests, and usage relative to requests
\citep{k8s_node_pressure_eviction}. The Actual Usage Evictor is not a
reclamation mechanism for an unhealthy Node. It protects a busy candidate from
a voluntary, strategy-driven Descheduler eviction.

\subsection{Research Gap}

Prior work demonstrates that runtime telemetry can define or repair placement
through scoring, candidate selection, destination selection, and migration.
The architectural gap addressed here is a strategy-independent, global
pre-eviction veto that leaves candidate ordering and placement objectives to
the host Descheduler strategy. The Actual Usage Evictor occupies this narrower
boundary by checking request-relative live usage immediately before eviction
and defining fail-closed behavior when the required runtime evidence is
unavailable or incomplete.
